\section{Positive time: the moving saddle and every fixed correction order}
\label{sec:positive}

\subsection{The saddle absorbs the large motion}
Write $\alpha=(\log M)'$. For each $w=s,1-s$, use the large
right-half-plane solution
\begin{equation}
 p=w+\frac t2\alpha(p),\qquad
 B=1-\frac t2\alpha'(p),\qquad
 \Sc_t(w)=M(p)e^{-(p-w)^2/t}B^{-1/2}.
 \label{eq:positive-saddle}
\end{equation}
The logarithm of $M$ is real on $p>1$ and $B^{-1/2}$ is the branch
near one. Fix the observation widths and a positive analytic margin;
all estimates start on the rectangle with twice that margin. For a
real center $x$ put
\begin{equation}
 X=|x|,\quad R=2+t+X,\quad \mathcal L_+=\log R,\quad
 \Ns=1+X+t\mathcal L_+,\qquad
 \A=|\Sc_t(s)|+|\Sc_t(1-s)|.
 \label{eq:positive-scales}
\end{equation}
At a sufficient onset, contraction in a fixed sector gives
\begin{equation}
 |p|\asymp\Ns,\quad \Re(p-w)\ge t\mathcal L_+/5,\quad
 |\Im(p-w)|\le t,\quad |B-1|\le1/4,\quad |p|\ge t.
 \label{eq:saddle-gates}
\end{equation}
These gates include bounded heights and the real center $x=0$.

\subsection{Finite coefficients}
Let $a_r=\alpha^{(r)}(p)$, $\beta_r=a_{r+1}/(r+2)!$, and
$\mathcal P_\ell=\{\nu\ge0:\sum r\nu_r=\ell\}$. With $k=|\nu|$,
define
\begin{align}
 A_\ell(u)&=\sum_{\nu\in\mathcal P_\ell}u^{\ell+2k}
                       \prod_r\frac{\beta_r^{\nu_r}}{\nu_r!},
 &\ell A_\ell&=\sum_{r=1}^\ell r\beta_r u^{r+2}A_{\ell-r},\\
 C_{2j}^{+}&=\sum_{\nu\in\mathcal P_{2j}}
 \frac{(2j+2k-1)!!}{2^{j+k}}
 \left(\frac tB\right)^{j+k}
 \prod_r\frac{a_{r+1}^{\nu_r}}{(r+2)!^{\nu_r}\nu_r!}.
 \label{eq:positive-coefficients}
\end{align}
The empty partition gives $C_0^+=1$. The superscript distinguishes
these coefficients from the negative-time Riemann--Siegel coefficients.
Parity kills every odd grade in the Gaussian moment. The first two
nonconstant coefficients are
\begin{align}
 C_2^+&=\frac{t^2\alpha'''(p)}{32B^2}
               +\frac{5t^3\alpha''(p)^2}{192B^3},\label{eq:c2}\\
 C_4^+&=\frac{t^3\alpha^{(5)}(p)}{384B^3}
 +\frac{t^4}{B^4}\left\{\frac{7\alpha''(p)\alpha^{(4)}(p)}{768}
                  +\frac{35\alpha'''(p)^2}{6144}\right\}\notag\\
 &\hspace{1em}+\frac{35t^5\alpha''(p)^2\alpha'''(p)}{2048B^5}
                  +\frac{385t^6\alpha''(p)^4}{73728B^6}.
 \label{eq:c4}
\end{align}
The approximation is
\begin{equation}
 P_{2J}(t,z)=\sum_{w=s,1-s}\Sc_t(w)\sum_{j=0}^J C_{2j}^+(t,w).
 \label{eq:positive-model}
\end{equation}
Every sum is finite. Nothing in this definition asserts convergence
of the infinite correction series.

\subsection{The complete error, before simplification}
For $n=2J+2$ define
\begin{equation}
 b_J=\sum_{k=1}^{2J+2}\frac{t^{J+1+k}}{\Ns^{2J+2+k}},\qquad
 \delta_J=b_J/t,\qquad
 E_J=b_J+R^{-\kappa_a t}+R^{b_0}e^{-ct\mathcal L_+^2},
 \quad\kappa_a=\frac{\log2}{8}.
 \label{eq:positive-budget}
\end{equation}
Here $b_0$ is a fixed middle-strip exponent and $c>0$ may be decreased
to at most $1/400$. The last two terms pay the complete arithmetic
remainder and the nonlocal contour pieces.

\begin{theorem}[Every fixed positive order]
\label{thm:positive-fixed}
Fix $J\ge0$ and fixed observation buffers. There are finite constants
and $R_{0,J}$ such that for $t\ge32$ and $R\ge R_{0,J}$,
\begin{align}
 |\partial_z^d(H-P_{2J})|&\le C_{J,d}\A\mathcal L_+^dE_J,
                         &&0\le d\le3,\\
 |(H-P_{2J})_t|&\le C_{t,J}\A(\mathcal L_+^2E_J+\delta_J),\\
 |\Res P_{2J}|&\le C_{r,J}\A\delta_J,
 &\Res&=\partial_t+\partial_{zz}.
 \label{eq:positive-full-error}
\end{align}
There is a sufficient $t_{*,J}$ for which, at every real height,
\begin{align}
 \frac{|\partial_z^d(H-P_{2J})|}{\A}
  &\le\frac{C_{J,d}}{t^{J+1}\log^{2J+3-d}(2+t)},&&0\le d\le3,\\
 \frac{|(H-P_{2J})_t|}{\A}
  &\le\frac{C_{t,J}}{t^{J+1}\log^{2J+1}(2+t)},\\
 \frac{|\Res P_{2J}|}{\A}
  &\le\frac{C_{r,J}}{t^{J+2}\log^{2J+3}(2+t)}.
 \label{eq:positive-eventual}
\end{align}
The field bound is first proved on the larger analytic rectangle;
the jet bounds use its smaller interior.
\end{theorem}
The full proof is now integrated in the ansatz companion's
\src{ansatz/manuscript/07_positive_hierarchy.tex}{\texttt{thm:positive-fixed}},
from the original Ansatz Note 19 \cite{positive}. Its
quantifier is ``for every fixed $J$, there exist constants and an
onset.'' The constants are specified by finite constructions and named
source constants, but are not numerically instantiated. The theorem
does not make $t=32$ an all-height onset for every $J$.

\subsection{Why the finite remainder has this form}
The following argument is the central reason the hierarchy is
systematically improvable at fixed order. Set $f=\log M$,
$a=\Re(p-w)$, $m=|p|$, and use the actual symmetric central segment
$\mathcal U=[-a/2,a/2]$. For $0\le\lambda\le1$ introduce
\begin{equation}
 E_\lambda(u)=u^2\int_0^1(1-r)
                  [f''(p+r\lambda u)-f''(p)]\dd r.
 \label{eq:interpolation}
\end{equation}
At $\lambda=1$ this is the exact nonquadratic exponent; at zero it
vanishes. Sectorial gamma estimates and Cauchy disks supply finite
$K_r$ with $|\alpha^{(r)}(q)|\le K_r|q|^{-r}$. After increasing the
onset, $|p+r\lambda u|\ge m/3$ and
\begin{equation}
 |\partial_\lambda^rE_\lambda|
 \le d_r\frac{|u|^{r+2}}{m^{r+1}},\qquad
 d_r=\frac{3^{r+1}K_{r+1}}{(r+1)(r+2)},\quad1\le r\le n.
 \label{eq:lambda-derivatives}
\end{equation}
The full interpolated Gaussian, not the nonquadratic factor by itself,
satisfies
\begin{equation}
 |e^{-Bu^2/t+E_\lambda(u)}|\le e^{-u^2/(2t)}.
 \label{eq:common-gaussian}
\end{equation}
Indeed its exponent is $-u^2/t+u^2\int_0^1(1-r)f''(p+r\lambda u)\dd r$;
the sufficient gate $3K_1t/(2m)\le1/2$ pays the second term.

Define the finite Bell constants
\[
 D_{n,k}=\sum_{\substack{\nu\in\mathcal P_n\\|\nu|=k}}
                   \prod_{r=1}^n\frac{(d_r/r!)^{\nu_r}}{\nu_r!}.
\]
Real Taylor's theorem in $\lambda$ through degree $n-1$ gives
\begin{equation}
 \left|e^{-Bu^2/t}\left(e^{E_1(u)}-\sum_{\ell=0}^{n-1}A_\ell(u)\right)\right|
 \le e^{-u^2/(2t)}\sum_{k=1}^nD_{n,k}
                            \frac{|u|^{n+2k}}{m^{n+k}}.
 \label{eq:bell-remainder}
\end{equation}
Every operation is finite, and only derivatives through
$\alpha^{(2J+3)}$ are needed. Integrating the even powers against
the Gaussian yields $t^{n/2+k}/m^{n+k}$, hence $b_J$.
Odd grades cancel on the actual symmetric segment. Only the retained
even polynomials are extended to whole Gaussian moments; their finite
tails are at most $C_J e^{-t\mathcal L_+^2/400}$.

The remainder proof is incomplete until the rest of the canonical
integral is paid. In \eqref{eq:positive-heat}, split at real parts
$2$ and $-1$, reflect the left arm using $\xi(w)=\xi(1-w)$, and
keep the original bounded middle. Shift the two right rays to their
saddles with their oriented connectors. Each shifted ray and connector
stays in a half-plane where the full ordinary Dirichlet series
converges. Both actual infinite ends, the noncentral parts, connectors,
middle, and all $n\ge2$ are retained in the error account. A gamma
integrand is never extended through its poles to an invented whole-line
contour. These payments give the last two terms of
\eqref{eq:positive-budget}. A sufficient eventual absorption includes
$\kappa_a t\ge2J+4$ and an order-dependent logarithmic onset.
This is why the full bound must precede its elegant simplified form.

\subsection{The residual recurrence and the direct time estimate}
The approximation has a nonzero heat residual, with a simple exact
description. In the saddle coordinate $p$ define
\begin{equation}
 \mathcal Q_+=\frac1{4B^2}\partial_{pp}
 +\frac{t\alpha''}{4B^3}\partial_p
 +\frac{t\alpha'''}{16B^3}
 +\frac{5t^2(\alpha'')^2}{64B^4}.
 \label{eq:positive-operator}
\end{equation}
Exact differentiation of \eqref{eq:positive-saddle} gives
\begin{equation}
 \Sc^{-1}(\partial_t-\tfrac14\partial_{ww})(\Sc F)
             =(\partial_t|_p-\mathcal Q_+)F,
 \qquad
 \partial_t|_p C_{2j+2}^+=\mathcal Q_+C_{2j}^+.
 \label{eq:positive-transport}
\end{equation}
One proves the coefficient identity with finite Gaussian moments and
$\rho=t/B$: differentiation uses $\rho_t|_p=B^{-2}$ and
$\rho_p=\alpha''\rho^2/2$. Thus
\begin{equation}
 \Res P_{2J}=-\sum_{w=s,1-s}\Sc_t(w)\mathcal Q_+C_{2J}^+,
 \qquad
 (H-P_{2J})_t=-(H-P_{2J})_{zz}-\Res P_{2J}.
 \label{eq:residual-telescope}
\end{equation}
The coefficient and transport proofs are in the ansatz companion's
\src{ansatz/manuscript/07_positive_hierarchy.tex}{\texttt{prop:positive-coefficients}, \texttt{prop:positive-residual}}.
The final grade is not identically zero; in particular it is eventually
nonzero at $z=0$. Formula~\eqref{eq:residual-telescope} supplies the
time estimate without differentiating a big-$O$ bound.

\subsection{The improvement has two useful faces}
For $X\gg t\mathcal L_+$, the allowance is
$b_J\asymp_J t^{J+2}X^{-2J-3}$. In the time-dominated or comparable
region it is $O_J(t^{-J-1}\mathcal L_+^{-2J-3})$. The next allowance satisfies
\begin{equation}
 \frac{b_{J+1}}{b_J}=O_J(t/\Ns^2)
              =O_J(1/(t\log^2(2+t))).
 \label{eq:order-gain}
\end{equation}
This compares proved positive bounds, not successive signed errors.
At each fixed height, a single arm's coefficient has asymptotic form
\[
 C_{2j}^+\sim\frac{2^{2j-1}(2j)!}{(j+1)!}
                     t^{-j}\log^{-2j-1}t\quad(j\ge1).
\]
The first omitted correction explains the field scales
$t^{-1}\log^{-3}t$, $t^{-2}\log^{-5}t$, and $t^{-3}\log^{-7}t$
for $P_0,P_2,P_4$. These sharpen the field description; existing
coarser physical-transfer theorems do not improve automatically.

\subsection{A compact-positive refinement reaches the switch}
Fix $I=[a,T]\subset[32,\infty)$, with finite $T$, and let $X\to\infty$.
The actual saddle equation gives
\[
 p-w=O_T(\log X),\quad |p|=X/2+O_T(\log X),\quad
 \Re(p-w)=\frac t4\log(X/(4\pi))+O_T(\log X/X).
\]
Using this position instead of the coarser all-height gate improves
the complete arithmetic allowance to $C_I X^{-t\log2/4}$.

\begin{theorem}[Compact positive field bridge]
\label{thm:compact-positive}
For each fixed $J,I$ and fixed buffers,
\begin{align}
 |\partial_z^d(H-P_{2J})|&\le C_{J,I,d}\A(1+\log X)^d
                      [X^{-2J-3}+X^{-t\log2/4}],&&0\le d\le3,\\
 |(H-P_{2J})_t|&\le C_{J,I}\A(1+\log X)^2
                      [X^{-2J-3}+X^{-t\log2/4}],\\
 |\Res P_{2J}|&\le C_{J,I}\A X^{-2J-3}.
 \label{eq:compact-positive}
\end{align}
In particular $P_2$ has pure $X^{-5}$ field allowance on every fixed
$[32,T]$, and $P_4$ has pure $X^{-7}$ allowance on $[63,65]$, with
the displayed jet and time losses.
\end{theorem}
The numerical inequalities are $8\log2>5$ and $63\log2/4>7$.
This is a new compact-domain proof, not substitution of $64$ into an
uninstantiated eventual onset. At higher orders the arithmetic term
must remain whenever its exponent does not pay the local precision.
It is an upper-bound budget, not a proved lower bound or a sharpness
claim. The full compact proof, finite-endpoint forcing and optional
splice are in the ansatz companion's
\src{ansatz/manuscript/07_positive_hierarchy.tex}{\texttt{thm:compact-positive}},
\src{ansatz/manuscript/07_positive_hierarchy.tex}{\texttt{prop:elementary-compact}}
and \src{ansatz/manuscript/07_positive_hierarchy.tex}{\texttt{prop:positive-overlap-splice}},
with original source Note 20 \cite{bridge}.

For comparison, an elementary corrected carrier is
\begin{equation}
 F_1=\sum_\pm C_\pm(1+D_\pm),\quad
 D_\pm=\frac{\pm\ii}{z}
  \left\{\frac1{24}-\frac t4-\frac{7t}4b_\pm-\frac{t^2}8b_\pm^2\right\},
 \quad b_\pm=L_z/2\pm\ii\pi/4.
 \label{eq:elementary-correction}
\end{equation}
With $\Ls=1+\log X$, the new complete comparison on $[3/4,33]$
is
\begin{align}
 B_1(t,X)&=\Ls^4/X^2+X^{-t\log2/4},\notag\\
 |\partial_z^j(H-F_1)|&\le C_j\Eo\Ls^jB_1,&&j\le3,\notag\\
 |(H-F_1)_t|&\le C_t\Eo\Ls^2B_1,\label{eq:intermediate-band}\\
 |\Res F_1|&\le C\Eo\Ls^4/X^2,
 &|\partial_z\Res F_1|&\le C\Eo\Ls^5/X^2.\notag
\end{align}
Its proof uses the complete fixed-endpoint canonical split and a
corrected local gamma calculation, retaining all arithmetic and the
actual middle and infinite ends \cite{intermediate}. It is integrated
as the ansatz companion's
\src{ansatz/manuscript/08_positive_bridges.tex}{\texttt{thm:intermediate-positive}}.
On $[3/4,1]$, comparison of $F_0$ and $F_1$ through $H$ pays the sum
of their complete allowances; comparison of $F_N$ and $F_1$ pays
$\epsilon_X+B_1$ with $\Ls^j$ for jets and $\Ls^2$ for time.
The omitted-arithmetic term in $B_1$ is retained. On $[32,33]$ this
is precisely the existing elementary $F_1$ formula; the common buffered
jet/time premises are checked, and $\A\asymp\Eo$. That identifies
its overlap with the compact fixed-order hierarchy without a new splice.
The truncated gamma arm's finite-endpoint forcing is separately paid;
it is not treated as an exact heat solution.

For a time profile $\chi$, any optional splice obeys
\begin{equation}
 \Res[\chi P_{2J}+(1-\chi)F_1]
 =\chi\Res P_{2J}+(1-\chi)\Res F_1+\chi'(P_{2J}-F_1).
 \label{eq:positive-splice}
\end{equation}
It inherits the coarser carrier error unless an additional cancellation
is proved. Choosing fixed $T>\max(65,t_{*,J})$ gives a high-height
overlap with the all-height theorem. The remaining bounded-height
compact core is still a separate approximation task.
